\documentclass[10.5pt,letterpaper]{article}
\usepackage[margin=2.0cm,top=1.6cm,bottom=1.8cm,headheight=14pt]{geometry}
\usepackage{fontspec}
\setmainfont{Arial}
\usepackage[spanish,es-noshorthands]{babel}
\usepackage{xcolor}
\usepackage{tabularx}
\usepackage{graphicx}
\usepackage{hyperref}
\usepackage{fancyhdr}
\usepackage{pgffor}
\usepackage{microtype}
\usepackage{setspace}

\definecolor{unabblue}{HTML}{173864}
\definecolor{linegray}{HTML}{BFBFBF}
\hypersetup{colorlinks=true,linkcolor=unabblue,urlcolor=unabblue}

\pagestyle{fancy}
\fancyhf{}
\fancyhead[L]{\scriptsize\sffamily\color{unabblue}UNIVERSIDAD DR. ANDRÉS BELLO - FACULTAD DE TECNOLOGÍA E INNOVACIÓN}
\fancyhead[R]{\scriptsize\sffamily\color{unabblue}Ciclo 02-2026}
\fancyfoot[L]{\scriptsize\sffamily\color{unabblue}Diseño de Sistemas de Seguridad de Redes de Datos}
\fancyfoot[R]{\scriptsize\sffamily\color{unabblue}Página \thepage\ de 3}
\renewcommand{\headrulewidth}{0.4pt}
\renewcommand{\footrulewidth}{0.4pt}

\singlespacing
\setlength{\parindent}{0pt}
\setlength{\parskip}{0pt}

% Renglones limpios estilo Word
\newcommand{\renglones}[1]{%
  \par\nopagebreak\vspace{0.4cm}%
  \foreach \i in {1,...,#1}{%
    \noindent{\color{linegray}\rule{\textwidth}{0.4pt}}\par\vspace{0.6cm}%
  }%
  \vspace{-0.15cm}%
}

\begin{document}

% =============================================================================
% PÁGINA 1
% =============================================================================
\begin{center}
  \includegraphics[width=1.6cm,height=1.6cm,keepaspectratio]{/home/jaimeberrios/Proyectos/UNAB/Ciclo-6/sDSR43/Plantilla-Documentos-Investigacion/assets/img/logo-portada.png}\\[0.08cm]
  {\normalsize\sffamily\bfseries\color{unabblue}UNIVERSIDAD DR. ANDRÉS BELLO}\\[0.02cm]
  {\small\sffamily\bfseries\color{unabblue}FACULTAD DE TECNOLOGÍA E INNOVACIÓN}\\[0.02cm]
  {\footnotesize\sffamily\color{unabblue}CARRERA: INGENIERÍA EN SISTEMAS Y COMPUTACIÓN}\\[0.02cm]
  {\scriptsize\sffamily Asignatura: Diseño de Sistemas de Seguridad de Redes de Datos - Ciclo 02-2026}\\[0.18cm]
  {\normalsize\sffamily\bfseries\color{unabblue}CUESTIONARIO DE DIAGNÓSTICO TÉCNICO DE INFRAESTRUCTURA DE RED}\\[0.08cm]
  {\color{unabblue}\rule{0.6\textwidth}{0.8pt}}
\end{center}

\vspace{0.1cm}

\textbf{\color{unabblue}\footnotesize I. DATOS GENERALES DEL DIAGNÓSTICO}
\vspace{0.05cm}

\begin{table}[h!]
  \centering
  \scriptsize
  \renewcommand{\arraystretch}{1.25}
  \begin{tabularx}{\textwidth}{|>{\bfseries\color{unabblue}\raggedright\arraybackslash}p{4.2cm}|X|}
    \hline
    Unidad de análisis: & Red de datos de la Sede de Enlaces con la Educación (Santa Rosa de Lima). \\
    \hline
    Lugar y fecha de aplicación: & \\
    \hline
    Nombre del informante: & \\
    \hline
    Cargo o función técnica: & \\
    \hline
    Equipo investigador: & Genesis Saravia, Jaime Berrios, Jefferson Bonilla, Rosmeri González, Florentín Treminio. \\
    \hline
  \end{tabularx}
\end{table}

\vspace{0.04cm}
{\scriptsize\itshape\color{black!80}Objetivo del instrumento: Recabar información técnica directa sobre el equipamiento, estado de obsolescencia física, topología, segmentación lógica y controles de filtrado perimetral de la sede institucional.}

\vspace{0.25cm}

\textbf{\color{unabblue}\footnotesize II. BATERÍA DE PREGUNTAS TÉCNICAS}

\vspace{0.15cm}

\subsection*{\color{unabblue}\footnotesize DIMENSIÓN 1: INFRAESTRUCTURA FÍSICA Y EQUIPAMIENTO CENTRAL}

\vspace{0.1cm}
\textbf{1. ¿Cuántas estaciones de trabajo operan activamente en la sede y bajo qué modalidades (cableada / inalámbrica) se distribuye la conectividad hacia los puestos del personal?}
\renglones{4}

\vspace{0.25cm}
\textbf{2. ¿Cuál es el dispositivo principal que concentra la distribución de red en el edificio, cuál es su tiempo aproximado de operación continua y en qué estado físico se encuentran sus puertos?}
\renglones{4}

\newpage

% =============================================================================
% PÁGINA 2
% =============================================================================

\textbf{3. ¿Cómo se encuentran administradas y distribuidas las salidas físicas del módem del proveedor de servicios de Internet?}
\renglones{4}

\vspace{0.5cm}

\subsection*{\color{unabblue}\footnotesize DIMENSIÓN 2: TOPOLOGÍA Y SEGMENTACIÓN LÓGICA DE RED}

\vspace{0.15cm}
\textbf{4. ¿Qué topología física rige la comunicación interna de la sede y existe actualmente alguna división lógica (VLAN) para separar los flujos de tráfico administrativo y técnico?}
\renglones{4}

\vspace{0.4cm}
\textbf{5. ¿Cuál es el tratamiento y canal de red utilizado cuando ingresan laptops externas de estudiantes o docentes para soporte técnico o diagnóstico?}
\renglones{4}

\newpage

% =============================================================================
% PÁGINA 3
% =============================================================================

\subsection*{\color{unabblue}\footnotesize DIMENSIÓN 3: SEGURIDAD PERIMETRAL Y CONTROLES DE TRÁFICO}

\vspace{0.15cm}
\textbf{6. ¿Existe actualmente algún servidor central o proxy activo encargado de la auditoría y el filtrado del tráfico hacia Internet, o qué antecedentes existen sobre dichos controles en la sede?}
\renglones{5}

\vspace{0.6cm}
\textbf{7. ¿Se aplican listas de control de acceso (ACL), políticas de puertos/servicios o mecanismos de inspección criptográfica (HTTPS/TLS) en los enlaces de red?}
\renglones{5}

\end{document}
